% Proposed B.89: actual discrete and compact symmetric filling comparison.
\statement{B.89}{Proved}{The actual symmetric filling intersection is the mutual commutator, with an exact compact fiber-product comparison}
\label{stmt:B89}
Use the actual surface and involution of B.87, and put
\[
 X=Y_b,\quad U=\mu _5\setminus\{1\},\quad
 A=\bigcup_{u\in U}\bigl(\{v=u\}\cup\{w=u\}\bigr)\cap X,
 \quad B=F(A).
\]
Let $Y=X\setminus A$, $V=X\setminus B=F(Y)$, and
$Z=Y\cap V$, all based at $P=(399/200,1/200)$. Write
\[
 \Gamma=\pi _1(Z,P),\quad G_1=\pi _1(Y,P),\quad
 G_2=\pi _1(V,P),\quad Q=\pi _1(X,P),
\]
and let $N_A=\ker(\Gamma\to G_2)$ and
$N_B=\ker(\Gamma\to G_1)$ be the actual filling kernels.
Then the actual maps induce
\begin{equation}\label{eq:B89-discrete}
 N_A\cap N_B=[N_A,N_B],\qquad
 \Gamma/[N_A,N_B]\xrightarrow{\sim}G_1\times_QG_2.
\end{equation}
Let $\mathcal N_A,\mathcal N_B$ be the corresponding closed filling
kernels in the full completion $\widehat\Gamma$. With brackets on
closed subgroups denoting the closed subgroup generated by their
commutators, one has
\begin{equation}\label{eq:B89-compact}
 \mathcal N_A\cap\mathcal N_B
   =[\mathcal N_A,\mathcal N_B]
   =\overline{N_A\cap N_B},
 \qquad
 \widehat{G_1\times_QG_2}
   \xrightarrow{\sim}\widehat G_1\times_{\widehat Q}\widehat G_2.
\end{equation}
The second isomorphism compares the actual groups and their full
completions. The comparison from $\Gamma$ itself need not be injective.
Neither equality asserts fixed-element lifting for the involution or
the original marginal source-kernel equality.

\begin{proof}
\emph{B.89.1. The actual open cover and its second homotopy group.}
The product map $q=vw$ makes $X$ a bundle over
$\mathbb C^*\setminus\mu _5$ with six-punctured-sphere fiber, as in
B.87. For $Y$ the fiber is
\[
 \mathbb P^1\setminus
  \bigl(\{0,\infty,\sqrt q,-\sqrt q\}\cup\mu _5
                       \cup q\mu _5\bigr).
\]
These fourteen points are distinct for $q\notin\{0\}\cup\mu _5$.
Local labeling and isotopy extension give the actual bundle, and the
same conclusion for $V$ follows by $F$. The base and fibers are
aspherical, so $X,Y,V$ are aspherical. Their homotopy exact sequences
give, in particular,
\begin{equation}\label{eq:B89-free-by-free}
 1\longrightarrow F_5\longrightarrow Q\longrightarrow F_6
       \longrightarrow1.
\end{equation}
Lifting the free base generators splits this group extension. No
geometric section is being asserted.

The intersection $A\cap B$ is finite. For example, for $u,t\in U$
the two component types containing $v=u$ have intersection coordinates
\[
 (v,w)=\left(u,\frac{u+t-ut}{u}\right),\qquad
 (v,w)=\left(u,\frac{t}{u+t-1}\right).
\]
The other types follow by exchanging $v,w$. The displayed divisors
have no common component; the denominator $u+t-1$ cannot vanish,
since two unit complex numbers summing to $1$ have quotient of order
three, whereas $u/t$ has order dividing five. The actual exact
$\mathbb Q(\zeta _5)$ certificate described below checks the retained
boundary conditions and finds all $64$ transverse mixed nodes. Only
the finiteness of this set $D=A\cap B$ is needed here.

Thus $Y\cup V=X\setminus D$. Removing a finite set from a smooth
real four-manifold induces isomorphisms on $\pi _1$ and $\pi _2$: relative
transversality moves maps of $S^2$, and their homotopies of dimension
at most three, off the finite set. In particular every nullhomotopy
in $X$ of a sphere in $X\setminus D$ can be chosen to avoid $D$.
It follows that
\begin{equation}\label{eq:B89-pi2}
 \pi _2(Y\cup V)=\pi _2(X)=0.
\end{equation}

\emph{B.89.2. The actual crossed-module calculation.}
The open cover $Y,V$ of $Y\cup V$ gives the homotopy pushout of
$Z\to Y$ and $Z\to V$. All spaces are connected. Each map on
fundamental groups is onto by moving loops off the added complex
divisor. The homotopy fibers are therefore connected. Because
$\pi _2(Y)=\pi _2(V)=0$, their homotopy exact sequences show both
that the maps
\[
 \pi _2\operatorname{hofib}(Z\to Y)\longrightarrow\pi _2(Z),
 \qquad
 \pi _2\operatorname{hofib}(Z\to V)\longrightarrow\pi _2(Z)
\]
are onto, and that the first homotopy groups of these fibers identify
with the inclusions of $N_B,N_A$ into $\Gamma$.

These are precisely the hypotheses of Lemma~4.1 of
\citet[p.~320]{BBrownLoday1987}. Its conclusion for the actual
homotopy pushout is
\begin{equation}\label{eq:B89-crossed}
 \pi _2(Y\cup V)\cong (N_A\cap N_B)/[N_A,N_B].
\end{equation}
Indeed the fiber product of the two inclusion crossed modules is
$N_A\cap N_B$, and their crossed commutator is the ordinary mutual
commutator; this identification is also explicit in
\citet[Propositions~2.5 and~2.8]{BBrown1984}. No asphericity of $Z$
is assumed or used. Equation~\eqref{eq:B89-pi2} proves the first
part of \eqref{eq:B89-discrete}.

Ordinary van Kampen and the codimension-four comparison give
$Q=\Gamma/(N_A N_B)$. For any two normal subgroups, the map to
the fiber product of the two quotient groups over their common
quotient is onto and has their intersection as kernel. To see
surjectivity, compatible representatives differ by a product from
$N_A N_B$; changing the representatives by its two factors makes
them equal. This proves the second part of
\eqref{eq:B89-discrete}.

\emph{B.89.3. Finite extensions of the actual filled group.}
We prove the following elementary fact. If
\[
 1\longrightarrow L\longrightarrow Q'\longrightarrow T
       \longrightarrow1
\]
has finitely generated free $L,T$, then every finite-kernel extension
$1\to D'\to E'\to Q'\to1$ has a finite-index subgroup meeting
$D'$ trivially. Taking its core yields a finite quotient of $E'$
which is injective on $D'$.

First suppose that $D'=C$ is central, and choose a splitting
$Q'=L\rtimes T$. Freeness gives homomorphic lifts
$i:L\to E'$ and $j:T\to E'$. For each $t\in T$ there is a
homomorphism $d_t:L\to C$ such that
\begin{equation}\label{eq:B89-mismatch}
 j(t)i(x)j(t)^{-1}=i(t\cdot x)d_t(x).
\end{equation}
The discrepancy is a homomorphism because $C$ is central. If $e$
is its exponent, the characteristic finite-index subgroup
$L_e=[L,L]L^e$ kills every $d_t$. Hence $j(T)$ normalizes $i(L_e)$,
and $i(L_e)j(T)$ maps isomorphically to $L_e\rtimes T$. It has
finite index in $E'$ and intersects $C$ trivially.

For arbitrary finite $D'$, the centralizer $C_{E'}(D')$ has finite
index. Its kernel over its image $Q'_0$ in $Q'$ is $Z(D')$, which
is central. The group $Q'_0$ has finite index in $Q'$; its intersection
with $L$ and its image in $T$ are finite-index finitely generated
free groups. The resulting extension splits because its quotient
is free. The central case applies, and taking the core in $E'$
proves the fact. In particular it applies to the actual $Q$ in
\eqref{eq:B89-free-by-free}, and to every finite-index subgroup
of that $Q$.

\emph{B.89.4. The full topology of a fiber product over this head.}
Let $q_i:G_i\twoheadrightarrow Q$ be the actual epimorphisms and
write $D_0=G_1\times_QG_2$, $R_i=\ker(q_i)$. Regard $R_i$ as
its coordinate subgroup in $D_0$. Let $K\triangleleft D_0$ have
finite index. The projection $D_0\to G_i$ is onto, so
$K\cap R_i$ is normal in $G_i$. The extension
\[
 1\longrightarrow R_i/(K\cap R_i)
  \longrightarrow G_i/(K\cap R_i)
  \longrightarrow Q\longrightarrow1
\]
has finite kernel. By B.89.3 there is a normal finite-index subgroup
$A_i\triangleleft G_i$ such that
\begin{equation}\label{eq:B89-coordinate-separation}
 A_i\cap R_i\subseteq K\cap R_i.
\end{equation}
Let $D_1=D_0\cap(A_1\times A_2)$ and let $Q_1$ be its image
in $Q$. These are finite-index subgroups. The kernel of
$D_1\to Q_1$ equals
$(A_1\cap R_1)\times(A_2\cap R_2)$ and is contained in $K$.
Consequently the finite homomorphism $D_1\to D_0/K$ factors
through $Q_1$. Take the core in $Q$ of the kernel of that finite
homomorphism on $Q_1$, and denote it by $Q_2$. It has finite index.
Then
\[
 D_0\cap\bigl((A_1\cap q_1^{-1}(Q_2))
                    \times(A_2\cap q_2^{-1}(Q_2))\bigr)\subseteq K.
\]
This proves that the full profinite topology on $D_0$ is induced
by the product of the full profinite topologies on $G_1,G_2$.

We also verify the image after completion. Consider a rectangular
open neighborhood of a point of
$\widehat G_1\times_{\widehat Q}\widehat G_2$, represented by
cosets $g_i A_i$ of normal finite-index subgroups. Compatibility
in $\widehat Q$ says that the cosets
$q_i(g_i)q_i(A_i)$ meet there. Each $q_i(A_i)$ has finite index
in $Q$, so their cosets already meet in $Q$: this is the same
finite-quotient condition modulo their intersection. Adjusting
the $g_i$ within $A_i$ gives an actual pair in $D_0$ in the given
neighborhood. Thus $D_0$ is dense in that compact fiber product.
Together with the induced topology, this proves the second
isomorphism of \eqref{eq:B89-compact}. This argument does not
invoke profinite homotopy excision or interchange intersections
and closures.

\emph{B.89.5. The actual closed intersection.}
The kernel of a completed surjective quotient is the closure of
its original normal kernel. Hence
$\mathcal N_A=\overline{N_A}$ and
$\mathcal N_B=\overline{N_B}$ inside $\widehat\Gamma$.
Density and continuity of commutation also give
\[
 \overline{[N_A,N_B]}=[\mathcal N_A,\mathcal N_B].
\]
Now complete \eqref{eq:B89-discrete} and use B.89.4. The resulting
isomorphisms are induced by the actual two filling maps:
\[
 \widehat\Gamma/\overline{[N_A,N_B]}
   \cong\widehat{G_1\times_QG_2}
   \cong\widehat G_1\times_{\widehat Q}\widehat G_2.
\]
The kernel of their composite from $\widehat\Gamma$ is also
$\mathcal N_A\cap\mathcal N_B$, which proves all of
\eqref{eq:B89-compact}.

\emph{B.89.6. Involution scope and exact evidence.}
The actual involution exchanges $N_A,N_B$ and their closures. After
identifying $G_2$ with $G_1$ by $F$, it exchanges the two coordinates
in the actual fiber products above. For a fixed target and any lift
$x$, the element $x^{-1}F(x)$ is a nonabelian $C_2$-cocycle in the
intersection kernel. Replacing $x$ by $xm$ changes it to
$m^{-1}x^{-1}F(x)F(m)$. Thus fixed lifting still requires that
particular class in nonabelian $H^1(C_2,N_A\cap N_B)$, or its
continuous analogue, to be neutral. Expressing this kernel as a
mutual commutator does not prove neutrality. No equality of the
projected fixed subgroup with the source image is inferred.

\begingroup\raggedright
The exact retained-boundary geometry is recorded in the directory
\path{work_DE/next_marginal_source/}, in
\path{check_symmetric_deletion_geometry.py}, its genuine output, and
\path{symmetric_deletion_geometry.json}.
The full proof and primary-source audit are retained in
\path{work_A/symmetric_intersection_commutator_16h.md}.
The finite-extension and topology arguments are proofs for all
finite quotients; no enumeration of their sizes is asserted.\par
\endgroup
\end{proof}
